\section*{Capítulo \ref{ch:g10:layered-ocean} --- O oceano em camadas}

\begin{solution}{exo:g10:layered-ocean:1}
A \omterm{def:g10:layered-ocean:mixed-layer}{camada de mistura} é a camada de superfície do oceano na qual a
temperatura e a salinidade quase não mudam com a profundidade; a \omterm{def:g10:layered-ocean:thermocline}{termoclina}
é a camada abaixo dela na qual a temperatura cai depressa com a
profundidade. O vento e as ondas agitam a \omterm{def:g10:layered-ocean:mixed-layer}{camada de mistura} (ajudados no
inverno pelo resfriamento da superfície).
\end{solution}

\begin{solution}{exo:g10:layered-ocean:2}
A \omterm{def:g10:layered-ocean:upwelling}{ressurgência} é a subida de água fria de baixo da camada de superfície,
que substitui a água de superfície empurrada para longe pelo vento. A água
que sobe traz nutrientes minerais para a camada iluminada, onde o
fitoplâncton os usa para crescer; o fitoplâncton alimenta o restante da
cadeia alimentar.
\end{solution}

\begin{solution}{exo:g10:layered-ocean:3}
Pacífico tropical: cerca de \qty{29}{\degreeCelsius} na superfície e
\qty{4.5}{\degreeCelsius} a \qty{1000}{m}. Oceano Ártico: cerca de
\qty{-1.4}{\degreeCelsius} na superfície e perto de
\qty{0}{\degreeCelsius} a \qty{1000}{m}.
\end{solution}

\begin{solution}{exo:g10:layered-ocean:4}
A \qty{0}{\degreeCelsius} (\qty{1028.1}{kg/m^3} contra \qty{1023.3}{kg/m^3}
a \qty{25}{\degreeCelsius}). A água do mar: a \qty{4}{\degreeCelsius}, a água
pura tem uma densidade de \qty{1000.0}{kg/m^3}, e a água do mar de
salinidade 35, cerca de \qty{1027.8}{kg/m^3} (lida no gráfico).
\end{solution}

\begin{solution}{exo:g10:layered-ocean:5}
O flutuador deriva a \qty{1000}{m}, afunda até \qty{2000}{m}, depois sobe
até a superfície medindo a temperatura, a salinidade (a partir da
condutividade) e a pressão no caminho; na superfície, ele envia o seu perfil
por satélite e afunda de novo. Os flutuadores derivam com as correntes por
todo o oceano e funcionam durante anos, em todas as estações, longe das
rotas dos navios.
\end{solution}

\begin{solution}{exo:g10:layered-ocean:6}
$361.84\times10^6 \times 3.682 \approx \qty{1.33e9}{km^3}$, cerca de
\num{1.33} bilhão de quilômetros cúbicos. $100/3682 \approx 0.027$: menos de
\qty{3}{\percent} da profundidade média.
\end{solution}

\begin{solution}{exo:g10:layered-ocean:7}
$22.7 - 0.5 = \qty{22.2}{\degreeCelsius}$, atingida a \qty{15}{m}: a \omterm{def:g10:layered-ocean:mixed-layer}{camada
de mistura} tem cerca de \qty{15}{m} de profundidade. Entre \qty{20}{m} e
\qty{50}{m}: $(21.8 - 17.3)/30 = \qty{0.15}{\degreeCelsius}$ por metro.
\end{solution}

\begin{solution}{exo:g10:layered-ocean:8}
O sol aquece apenas a camada de superfície, e o aquecimento e o resfriamento
sazonais só se espalham pela \omterm{def:g10:layered-ocean:mixed-layer}{camada de mistura}: abaixo da \omterm{def:g10:layered-ocean:thermocline}{termoclina}, a água
fica isolada deles. No inverno, a superfície se resfria, fica mais densa e
afunda na água de baixo, e as tempestades agitam o mar: a mistura vai
fundo. No verão, o sol aquece uma camada de superfície leve que flutua
sobre a água mais fria; os ventos mais suaves não conseguem empurrá-la para
baixo, e a \omterm{def:g10:layered-ocean:mixed-layer}{camada de mistura} continua fina.
\end{solution}

\begin{solution}{exo:g10:layered-ocean:9}
$0.78/0.18 \approx 4.3$ graus perto de \qty{10}{\degreeCelsius};
$0.78/0.06 = 13$ graus perto de \qty{0}{\degreeCelsius}. Nos mares polares
frios, pequenas mudanças de salinidade pesam muito mais na densidade que as
mudanças de temperatura: ali, é a salinidade que controla as camadas.
\end{solution}

\begin{solution}{exo:g10:layered-ocean:10}
A temperatura é a mesma: a água mais funda é mais densa porque é mais
salgada (33.3 contra 28.1). A camada entre elas, onde a salinidade muda
depressa com a profundidade, é uma \omterm{def:g10:layered-ocean:thermocline}{haloclina}; aqui ela é também a
\omterm{def:g10:layered-ocean:thermocline}{picnoclina}.
\end{solution}

\begin{solution}{exo:g10:layered-ocean:11}
Empurrada para baixo, a parcela conserva a sua temperatura e a sua
salinidade, logo a sua densidade. Cercada de água mais densa, ela é mais
leve que a água que desloca: o empuxo a leva de volta ao seu próprio nível.
Uma parcela empurrada para cima seria mais densa que o seu novo entorno e
voltaria a afundar. As camadas são estáveis.
\end{solution}

\begin{solution}{exo:g10:layered-ocean:12}
Massa: $0.001 \times 1028.1 = \qty{1.0281}{kg}$. Novo volume: $1.0281/1023.3
\approx \qty{0.0010047}{m^3}$, cerca de \qty{1004.7}{mL}: ela se dilatou
cerca de \qty{4.7}{mL}, \qty{0.47}{\percent}. A camada de \qty{700}{m}
engrossaria $700 \times 0.0047 \approx \qty{3.3}{m}$.
\end{solution}

\begin{solution}{exo:g10:layered-ocean:13}
No Atlântico Norte em fevereiro. Ali, os \qty{150}{m} de cima têm quase a
mesma temperatura, logo quase a mesma densidade: misturá-los custa pouco.
No Pacífico tropical, a densidade sobe $1026.0 - 1021.6 =
\qty{4.4}{kg/m^3}$ entre \qty{50}{m} e \qty{200}{m}: uma \omterm{def:g10:layered-ocean:thermocline}{picnoclina} forte, e
a tempestade teria de empurrar água leve para dentro de uma água muito mais
densa.
\end{solution}

\begin{solution}{exo:g10:layered-ocean:14}
Clorofila alta longe das costas: a faixa ao longo do equador no Pacífico (e
no Atlântico), onde os ventos alísios provocam uma \omterm{def:g10:layered-ocean:upwelling}{ressurgência}; e as altas
latitudes, como o Atlântico Norte, onde a \omterm{def:g10:layered-ocean:mixed-layer}{camada de mistura} de inverno
chega fundo e traz os nutrientes para cima. Clorofila muito baixa: os
centros dos giros subtropicais, fortemente \omterm{def:g10:layered-ocean:stratification}{estratificados}, onde nada traz
água profunda para cima e a camada iluminada ficou sem nutrientes.
\end{solution}

\begin{solution}{exo:g10:layered-ocean:15}
Descida: $1000/(6 \times 3600) \approx \qty{0.046}{m/s}$, cerca de
\qty{4.6}{cm/s}. Subida: $2000/(10 \times 3600) \approx \qty{0.056}{m/s}$,
cerca de \qty{5.6}{cm/s}. $36\,000/2 = 18\,000$ medições, em média a
$2000/18\,000 \approx \qty{0.11}{m}$ uma da outra.
\end{solution}

\begin{solution}{pb:g10:layered-ocean:1}
\textbf{1.} O giro, por $24.0 - 19.0 = \qty{5.0}{\degreeCelsius}$.

\textbf{2.} O giro (\qty{20.5}{\degree S} contra \qty{13.5}{\degree S}).
Mais longe do equador, ele deveria ser o mais fresco: as latitudes sozinhas
não podem explicar por que ele é \qty{5}{\degreeCelsius} mais quente.

\textbf{3.} $24.0 - 0.5 = \qty{23.5}{\degreeCelsius}$, atingida entre
\qty{50}{m} (23.6) e \qty{100}{m} (22.6), a cerca de $50 + 50 \times
0.1/1.0 = \qty{55}{m}$.

\textbf{4.} Não: a temperatura cai regularmente desde a superfície,
\qty{1.6}{\degreeCelsius} nos primeiros \qty{25}{m} e
\qty{5.3}{\degreeCelsius} nos primeiros \qty{100}{m}; a \omterm{def:g10:layered-ocean:thermocline}{termoclina} começa na
superfície.

\textbf{5.} Entre \qty{200}{m} e \qty{300}{m}: $20.3 - 14.4 =
\qty{5.9}{\degreeCelsius}$ por cem metros.

\textbf{6.} Ao largo do Peru, entre \qty{50}{m} e \qty{100}{m}: cerca de
\qty{70}{m} por interpolação (o perfil completo dá \qty{65}{m}). No giro,
entre \qty{200}{m} e \qty{300}{m}: cerca de $200 + 100 \times 5.3/5.9
\approx \qty{290}{m}$.

\textbf{7.} Sim, em todos os níveis das duas colunas: água mais leve repousa
sobre água mais densa, e as duas colunas são estáveis.

\textbf{8.} A água mais salgada deveria ser mais densa; no entanto, a água
de superfície do giro é mais leve: \qty{1024.71}{kg/m^3} contra
\qty{1025.08}{kg/m^3}.

\textbf{9.} A água do giro é cerca de 1.3 unidade mais salgada, o que
acrescenta cerca de $1.3 \times 0.8 \approx \qty{1.0}{kg/m^3}$; ela é
\qty{5}{\degreeCelsius} mais quente, o que retira cerca de
$5 \times 0.27 \approx \qty{1.4}{kg/m^3}$ nessas temperaturas. A
temperatura vence.

\textbf{10.} $1026.44 - 1025.08 = \qty{1.36}{kg/m^3}$.

\textbf{11.} $1025.30 - 1024.71 = \qty{0.59}{kg/m^3}$.

\textbf{12.} Da superfície até \qty{500}{m}, a densidade sobe
$1026.92 - 1024.71 = \qty{2.21}{kg/m^3}$, dos quais apenas
\qty{0.59}{kg/m^3} acima de \qty{200}{m}: a maior parte,
\qty{1.62}{kg/m^3}, fica abaixo de \qty{200}{m}, na \omterm{def:g10:layered-ocean:thermocline}{termoclina} entre
\qty{200}{m} e \qty{500}{m}.

\textbf{13.} Ao largo do Peru, a densidade sobe $1026.97 - 1025.08 =
\qty{1.89}{kg/m^3}$ da superfície até \qty{500}{m}, dos quais
\qty{1.36}{kg/m^3} nos \qty{200}{m} de cima: a maior parte fica perto da
superfície.

\textbf{14.} A \qty{500}{m}, os dois locais são quase iguais
(\qty{7.9}{\degreeCelsius} e \qty{7.1}{\degreeCelsius},
\qty{1026.97}{kg/m^3} e \qty{1026.92}{kg/m^3}): a mesma água fria fica sob
os dois; a diferença é o quanto ela chega perto da superfície.

\textbf{15.} O vento que sopra ao longo da costa em direção ao equador
empurra a água de superfície para o mar aberto; a água de baixo sobe perto
da costa para tomar o seu lugar. A água fria da \omterm{def:g10:layered-ocean:thermocline}{termoclina} é puxada para
cima, de modo que a \omterm{def:g10:layered-ocean:thermocline}{termoclina} se inclina para cima em direção à costa e
chega quase até a superfície.

\textbf{16.} Ao largo do Peru: a água a \qty{13}{\degreeCelsius} a
\qty{15}{\degreeCelsius}, de baixo da \omterm{def:g10:layered-ocean:thermocline}{termoclina}, fica nos \qty{100}{m} de
cima, ao alcance da agitação pelo vento e da luz. No giro, ela fica abaixo
de \qty{250}{m}, sob uma camada quente e leve que o vento não consegue
agitar até essa profundidade.

\textbf{17.} A \omterm{def:g10:layered-ocean:upwelling}{ressurgência} pararia: a água de superfície se aqueceria ao
sol, uma camada quente e leve se formaria por cima, e o fornecimento de
nutrientes à camada iluminada diminuiria.

\textbf{18.} A costa do Peru vai do azul-claro ao amarelo (muita clorofila);
o meio do giro do Pacífico Sul está entre as regiões de azul mais escuro do
mapa (muito pouca).

\textbf{19.} O fitoplâncton precisa de nutrientes além de luz. No giro, a
forte \omterm{def:g10:layered-ocean:stratification}{estratificação} mantém a água profunda, rica em nutrientes, muito
abaixo da camada iluminada, e os nutrientes da camada de superfície já foram
consumidos há muito tempo.

\textbf{20.} $1.36/0.59 \approx 2.3$: \textbf{o aumento de densidade nos
\qty{200}{m} de cima é cerca de 2.3 vezes maior ao largo do Peru que no
giro}. A \omterm{def:g10:layered-ocean:thermocline}{picnoclina}, com a água fria e rica em nutrientes abaixo dela, foi
levantada para dentro dos \qty{200}{m} de cima ao largo do Peru, enquanto no
giro ela fica mais funda, abaixo de \qty{200}{m}.
\end{solution}
